% Generated draft: requires mathematical and source review.
\documentclass[11pt]{article}
\usepackage[margin=25mm]{geometry}
\usepackage{fontspec}
\setmainfont{Latin Modern Roman}
\usepackage{amsmath,amssymb,mathtools,unicode-math}
\setmathfont{Latin Modern Math}
\usepackage{enumitem,needspace,xurl,hyperref}
\setlength{\parindent}{0pt}
\setlength{\parskip}{5pt}
\setlength{\emergencystretch}{3em}
\newcommand{\N}{\mathbb{N}}
\newcommand{\Z}{\mathbb{Z}}
\newcommand{\Q}{\mathbb{Q}}
\newcommand{\R}{\mathbb{R}}
\newcommand{\C}{\mathbb{C}}
\newcommand{\F}{\mathbb{F}}
\newcommand{\A}{\mathbb{A}}
\newcommand{\PP}{\mathbb{P}}
\newcommand{\E}{\mathbb{E}}
\newcommand{\eps}{\varepsilon}
\newcommand{\dd}{\,\mathrm{d}}
\begin{document}
\section*{UnsolvedMath ID: 6700023}
\textbf{Source problem code:} catalogue entry
\subsection{Definitions and mathematical statement}
The source tentatively proposes, at least for compact $X$,
\[
\mathrm{K\!\text{-}area}(X\times\mathbb R)
=\mathrm{K\!\text{-}area}(X\times S^1).
\]
This is a metric-dependent statement. If $(X,g)$ is the intended compact Riemannian manifold and $S^1_\ell$ is a circle of circumference $\ell>0$, an explicit product-metric version would compare
\[
\mathrm{K\!\text{-}area}(X\times\mathbb R,g+dt^2)
\quad\text{and}\quad
\mathrm{K\!\text{-}area}(X\times S^1_\ell,g+ds^2).
\]
The source does not specify $\ell$ or whether the assertion concerns all circle sizes, sufficiently small sizes, or a limiting value as $\ell\to0$.

In a conventional curvature-based K-area construction, one takes a supremum of reciprocal curvature norms of admissible Hermitian bundles with unitary connections and nontrivial characteristic data. The admissibility, nontriviality, parity/stabilization, and conditions at infinity are essential to the definition. They are not supplied here, and the summary calls the invariant stable K-area without giving a stabilization convention. Consequently no particular bundle formula is substituted for the source's invariant. The intended equality should be formulated with one specified K-area variant, the same conventions on both factors, and explicit product metrics.
\subsection{Short English statement}
For compact $X$, compare K-area after adjoining a line with K-area after adjoining a circle. Fix the invariant variant, conditions at infinity, product metrics, and circle-size quantifier before asserting equality.
\subsection{Sources}
\begin{itemize}
\item Misha Gromov, \emph{101 Questions, Problems and Conjectures around Scalar Curvature} (2017), printed tag [?22], occurrence 23, PDF page 32, \url{https://www.ihes.fr/~gromov/wp-content/uploads/2018/08/101-problemsOct1-2017.pdf}; accessed 2026-07-29, PDF-text-normalized extraction. The provenance describes the public author PDF as incomplete.
\item Local registry \texttt{data/unsolvedmath/problems.json}, permanent UnsolvedMath ID $6700023$, source code catalogue entry, \emph{Scalar Curvature Question [?22]}. Dataset attribution: Ulam AI and the UnsolvedMath Contributors.
\end{itemize}
\end{document}
